\mxfront{terminale-e9}
\mxnewpage\phantomsection\label{sujet:9-1}\addcontentsline{toc}{chapter}{Blanc final · Zéro 1}

\section*{Blanc final · Zéro 1}

\mxbanner{SUJET ZÉRO · ENTRAÎNEMENT}

\textbf{Durée : 4 h · Note sur 20 · \mxTallied{ex-9-1} exercices indépendants.}\\Calculatrice autorisée. Toute réponse doit être justifiée ; les démarches incomplètes seront prises en compte. Donner les valeurs exactes ; arrondir les probabilités à $10^{-6}$ et les autres approximations à $10^{-4}$, sauf précision différente.

\smallskip\textbf{Nom :}\enspace\dotfill\quad\textbf{Classe :}\enspace\dotfill\quad\textbf{Date :}\enspace\dotfill

\mxtally{ex-9-1}
\subsection*{Exercice 1 · Une fonction logarithmique \hfill 5 points}On considère la fonction $f$ définie sur $\intoo{0}{+\infty}$ par $f(x)=\frac{2+\ln x}x$. On note $\mathcal C_f$ sa courbe représentative.

\begin{enumerate}

\item Déterminer les limites de $f$ en $0^+$ et en $+\infty$. \hfill\textit{(\num{0,5} pt)}

\item Calculer $f' (x)$, puis dresser le tableau de variations de $f$. \hfill\textit{(\num{0,5} pt)}

\item Préciser la valeur du maximum de $f$, puis résoudre $f(x)=0$. \hfill\textit{(\num{0,5} pt)}

\item Étudier la convexité de $f$ et déterminer les coordonnées du point d'inflexion de $\mathcal C_f$. \hfill\textit{(1 pt)}

\item Démontrer que l'équation $f(x)=\frac12 e^{1}$ admet exactement deux solutions dans $\intoo{0}{+\infty}$. \hfill\textit{(\num{0,5} pt)}

\item Déterminer une primitive de $f$ sur $\intoo{0}{+\infty}$. \hfill\textit{(\num{0,5} pt)}

\item Calculer la valeur moyenne de $f$ sur $\intff{1}{e}$. \hfill\textit{(\num{0,5} pt)}

\item Indépendamment des questions précédentes, résoudre $\ln(x-2)+\ln(x+2)=\ln(32)$, après avoir précisé l'ensemble sur lequel cette équation est définie. \hfill\textit{(1 pt)}

\end{enumerate}

\mxtally{ex-9-1}
\subsection*{Exercice 2 · Phénomène périodique \hfill 5 points}Un point se déplace sur une roue tournant à vitesse constante dans un plan vertical. Sa hauteur au-dessus du sol est modélisée par $h(t)=9-3\cos\left(\frac{\pi t}{5}\right)$, où $t$ est exprimé en secondes et $h(t)$ en mètres. On étudie un tour complet, pour $t\in\intff{0}{10}$.

\begin{enumerate}

\item Donner une période du modèle prolongé à $\R$, puis les hauteurs minimale et maximale sur un tour. \hfill\textit{(\num{0,5} pt)}

\item Calculer $h'(t)$ et dresser le tableau de variations de $h$ sur $\intff{0}{10}$. \hfill\textit{(\num{0,5} pt)}

\item Déterminer tous les instants $t\in\intff{0}{10}$ tels que $h(t)=\frac{21}{2}$. \hfill\textit{(\num{0,5} pt)}

\item Calculer la valeur moyenne de $h$ sur $\intff{0}{10}$. \hfill\textit{(1 pt)}

\item Étudier la convexité de $h$ sur l'intervalle d'étude et déterminer les coordonnées des points d'inflexion de sa courbe. \hfill\textit{(\num{0,5} pt)}

\item Déterminer la durée totale pendant laquelle la hauteur est au moins égale à $\frac{21}{2}$ m. \hfill\textit{(\num{0,5} pt)}

\item La vitesse verticale est donnée par $h'(t)$. Déterminer ses valeurs maximale et minimale, en mètres par seconde, ainsi que les instants où elles sont atteintes. \hfill\textit{(\num{0,5} pt)}

\item Calculer le déplacement vertical $h(10)-h(0)$, puis la somme de la hauteur parcourue à la montée et de celle parcourue à la descente. Expliquer pourquoi ces deux résultats sont différents. \hfill\textit{(1 pt)}

\end{enumerate}

\mxnewpage\mxcontinued{Blanc final · Zéro 1 · suite}

\mxtally{ex-9-1}
\subsection*{Exercice 3 · Droite, plan et projeté \hfill 5 points}Dans un repère orthonormé de l'espace, on considère le plan $\mathcal P$ d'équation $x+2y+2z=11$ et le point $M\pt{2;0;0}$. La droite $d$ passe par $M$ et admet pour vecteur directeur $\vec n=\pt{1;2;2}$.

\begin{enumerate}

\item Donner une représentation paramétrique de $d$. \hfill\textit{(\num{0,5} pt)}

\item Démontrer que $d$ coupe $\mathcal P$ en un point $H$, puis calculer les coordonnées de $H$. \hfill\textit{(\num{0,5} pt)}

\item Justifier que $H$ est le projeté orthogonal de $M$ sur $\mathcal P$, puis calculer la distance de $M$ à $\mathcal P$. \hfill\textit{(\num{0,5} pt)}

\item On considère le plan $\mathcal Q$ d'équation $y=z$. Démontrer que $\mathcal P\cap\mathcal Q$ est une droite et en donner une représentation paramétrique. \hfill\textit{(1 pt)}

\item Vérifier que $H\in\mathcal Q$. En déduire la distance de $M$ à la droite $\mathcal P\cap\mathcal Q$. \hfill\textit{(\num{0,5} pt)}

\item Déterminer une équation cartésienne du plan parallèle à $\mathcal P$ passant par $M$. \hfill\textit{(\num{0,5} pt)}

\item Déterminer l'angle entre $d$ et la droite $\mathcal P\cap\mathcal Q$. \hfill\textit{(\num{0,5} pt)}

\item On définit le point $N$ par $\vect{MN}=\vec n$. Calculer les distances de $N$ au plan $\mathcal P$ et à la droite $\mathcal P\cap\mathcal Q$, puis les comparer. \hfill\textit{(1 pt)}

\end{enumerate}

\mxtally{ex-9-1}
\subsection*{Exercice 4 · Répétitions et estimation \hfill 5 points}On répète $20$ fois une épreuve à deux issues, appelées succès et échec. Les épreuves sont indépendantes et la probabilité de succès est $p=\frac{1}{5}$ à chaque épreuve. On note $X$ le nombre de succès obtenus.

\begin{enumerate}

\item Justifier que $X$ suit une loi binomiale et préciser ses paramètres. \hfill\textit{(\num{0,5} pt)}

\item Calculer la probabilité d'obtenir exactement deux succès. \hfill\textit{(\num{0,5} pt)}

\item Exprimer à l'aide des événements contraires, puis calculer, la probabilité d'obtenir au moins deux succès. \hfill\textit{(\num{0,5} pt)}

\item Calculer l'espérance et la variance de $X$. \hfill\textit{(1 pt)}

\item On répète maintenant la même épreuve un nombre $N\ge1$ de fois, toujours indépendamment. Déterminer le plus petit entier $N$ pour lequel la probabilité d'obtenir au moins un succès est supérieure ou égale à $\num{0,95}$. Une recherche par tableau ou par boucle est autorisée. \hfill\textit{(\num{0,5} pt)}

\item Pour un entier $j\in\{0,\ldots,20\}$, démontrer la formule de $P(X=j)$ en dénombrant les positions possibles des succès. \hfill\textit{(\num{0,5} pt)}

\item Pour $N$ épreuves indépendantes de même paramètre $p$, on note $F_N$ la fréquence des succès. Démontrer les expressions de $E(F_N)$ et de $V(F_N)$ en fonction de $N$ et de $p$. \hfill\textit{(\num{0,5} pt)}

\item À l'aide de Bienaymé-Tchebychev, déterminer une taille $N$ suffisante pour que $P(|F_N-p|<\num{0,02})\ge\num{0,95}$. La taille obtenue est-elle nécessairement la plus petite possible pour cette probabilité ? Justifier. \hfill\textit{(1 pt)}

\end{enumerate}

\ifmxteacher\mxnewpage

\section*{Corrigé · Blanc final · Zéro 1}\label{cor:9-1}

\subsection*{Exercice 1 · Une fonction logarithmique}\begin{enumerate}

\item \textit{[\num{0,5} pt]} En $0^+$, le numérateur tend vers $-\infty$ et le dénominateur vers $0^+$ : limite $-\infty$ ; en $+\infty$, limite $0$ par croissance comparée.

\item \textit{[\num{0,5} pt]} Par la règle du quotient,
\[f'(x)=\frac{(\frac{1}{x})x-(2+\ln x)}{x^2}=\frac{-1-\ln x}{x^2}.\]
Le dénominateur est positif. Le numérateur s'annule en $s=e^{-1}$, est positif pour $0<x<s$ et négatif pour $x>s$. Ainsi $f$ croît strictement sur $]0,s]$, puis décroît strictement sur $[s,+\infty[$.

\item \textit{[\num{0,5} pt]} Le maximum est atteint uniquement en $s=e^{-1}$ et vaut
\[f(s)=\frac{2- 1}{e^{-1}}=e^{1}.\]
Comme $x>0$, $f(x)=0$ équivaut à $2+\ln x=0$, soit $x=e^{-2}$.

\item \textit{[1 pt]} On dérive $(-1-\ln x)x^{-2}$ :
\[f''(x)=-x^{-3}-2(-1-\ln x)x^{-3}=\frac{2\ln x+1}{x^3}.\]
Le numérateur s'annule en $t=e^{-\frac{1}{2}}$ et passe du signe négatif au signe positif. La fonction est concave sur $]0,t]$, convexe sur $[t,+\infty[$. Elle change de convexité en $t$ et $f(t)=\frac{3}{2t}$ : le point d'inflexion est $I\pt{t;\frac{3}{2t}}$.

\item \textit{[\num{0,5} pt]} Posons $m=\frac{e^{1}}{2}$ et $s=e^{-1}$. La fonction est continue sur $]0,+\infty[$. Sur $]0,s]$, elle est strictement croissante, de limite $-\infty$ en $0$ et de valeur $2m$ en $s$ : le corollaire du TVI donne exactement une solution dans $]0,s[$. Sur $[s,+\infty[$, elle est strictement décroissante de $2m$ vers $0<m$ : exactement une solution dans $]s,+\infty[$. Ces intervalles couvrent le domaine et $f(s)\ne m$ : il y a exactement deux solutions.

\item \textit{[\num{0,5} pt]} Posons $F(x)=2\ln x+\frac12(\ln x)^2$, dérivable pour $x>0$. Alors $F'(x)=\frac{2}{x}+\frac{\ln x}{x}=f(x)$ ; $F$ est une primitive sur $]0,+\infty[$.

\item \textit{[\num{0,5} pt]} Avec cette primitive, $F(e)=\frac{5}{2}$ et $F(1)=0$. L'intégrale vaut donc $\frac{5}{2}$, et la valeur moyenne est
\[\mu=\frac{1}{e-1}\int_1^e f(x)\,\dd x=\frac{5}{2(e-1)}.\]

\item \textit{[1 pt]} Domaine $x>2$ ; $x^2-4=32$, donc $x=6$, la racine négative étant exclue.

\end{enumerate}

\subsection*{Exercice 2 · Phénomène périodique}\begin{enumerate}

\item \textit{[\num{0,5} pt]} Période $10$ s ; minimum $6$ m, maximum $12$ m.

\item \textit{[\num{0,5} pt]} $h'(t)=(\frac{3}{5}\pi)\sin (\frac{\pi t}{5})$ ; croissance sur $[0,5]$, décroissance sur $[5,10]$.

\item \textit{[\num{0,5} pt]} L'équation équivaut à $\cos (\frac{\pi t}{5})=-\frac{1}{2}$. Sur un tour, $\frac{\pi t}{5}\in[0,2\pi]$ ; les angles convenables sont $\frac{2\pi}{3}$ et $\frac{4\pi}{3}$. Ainsi $t=\frac{10}{3}$ s ou $t=\frac{20}{3}$ s.

\item \textit{[1 pt]} Une primitive est $F(t)=9t-\frac{15}{\pi}\sin (\frac{\pi t}{5})$. Comme $\sin0=\sin(2\pi)=0$,
\[\overline h=\frac{1}{10}\int_0^{10}h(t)\,\dd t=\frac{F(10)-F(0)}{10}=\frac{90}{10}=9\ \mathrm{m}.\]

\item \textit{[\num{0,5} pt]} La dérivée seconde est $h''(t)=\frac{3\pi^2}{25}\cos (\frac{\pi t}{5})$. Son signe est celui du cosinus : positif sur $[0,\frac{5}{2}[$ et $]\frac{15}{2},10]$, négatif entre ces deux intervalles. La fonction est convexe sur $[0,\frac{5}{2}]$ et $[\frac{15}{2},10]$, concave sur $[\frac{5}{2},\frac{15}{2}]$. Les deux changements de convexité donnent les points d'inflexion $\pt{\frac{5}{2};9}$ et $\pt{\frac{15}{2};9}$.

\item \textit{[\num{0,5} pt]} Entre $\frac{10}{3}$ et $\frac{20}{3}$ ; durée $\frac{10}{3}$ s.

\item \textit{[\num{0,5} pt]} $h'$ atteint $\frac{3}{5}\pi$ $\mathrm{m\,s^{-1}}$ en $t=\frac{5}{2}$ et son opposé en $t=\frac{15}{2}$.

\item \textit{[1 pt]} Déplacement $h(10)-h(0)=0$. Montée de $6$ m puis descente de $6$ m : distance cumulée $12$ m, égale à l'intégrale de $|h'|$.

\end{enumerate}

\subsection*{Exercice 3 · Droite, plan et projeté}\begin{enumerate}

\item \textit{[\num{0,5} pt]} $\pt{x;y;z}=\pt{2+t;2t;2t}$, $t\in\R$.

\item \textit{[\num{0,5} pt]} $2+9t=11$ donne $t=1$ ; $H=\pt{3;2;2}$.

\item \textit{[\num{0,5} pt]} Le vecteur $\vec n=\pt{1;2;2}$ est normal à $\mathcal P$. Comme $H\in\mathcal P$ et $\vect{MH}=\vec n$, $H$ est le projeté orthogonal de $M$ sur le plan. Le repère étant orthonormé, $\|\vec n\|=\sqrt{1+4+4}=3$. Donc
\[d(M,\mathcal P)=MH=|1|\,3=3.\]

\item \textit{[1 pt]} Poser $y=z=s$ ; $x=11-4s$. Donc $\pt{x;y;z}=\pt{11;0;0}+s\pt{-4;1;1}$.

\item \textit{[\num{0,5} pt]} Oui, ses deux dernières coordonnées sont égales. $H$ est dans la droite et $\vect{MH}$ est orthogonal à toute direction du plan, donc à cette droite. Distance $3$.

\item \textit{[\num{0,5} pt]} $x+2y+2z=2$ ; même vecteur normal et passage par $M$.

\item \textit{[\num{0,5} pt]} $\vec n\cdot\pt{-4;1;1}=-4+2+2=0$ : angle $\frac{\pi}{2}$.

\item \textit{[1 pt]} $N=M+\vec n$ est sur la même droite normale que $M$ ; son projeté sur $\mathcal P$ est donc aussi $H$. Or $H\in\mathcal P\cap\mathcal Q$. Pour tout $P$ de cette droite, $NP^2=NH^2+HP^2$, donc sa distance minimale à $N$ est $NH$. Finalement,
\[d(N,\mathcal P)=d(N,\mathcal P\cap\mathcal Q)=\|\vect{NH}\|=|1-1|\,\|\vec n\|=0.\]

\end{enumerate}

\subsection*{Exercice 4 · Répétitions et estimation}\begin{enumerate}

\item \textit{[\num{0,5} pt]} Répétition d'un nombre fixe $20$ d'épreuves identiques indépendantes à deux issues ; $X\sim\mathcal B\pt{20;\frac{1}{5}}$.

\item \textit{[\num{0,5} pt]} $P(X=2)=\binom{20}2(\frac{1}{5})^2(\frac{4}{5})^{18}\simeq\num{0,136909}$.

\item \textit{[\num{0,5} pt]} $1-P(X=0)-P(X=1)=1-(\frac{4}{5})^{20}-20(\frac{1}{5})(\frac{4}{5})^{19}\simeq\num{0,930825}$.

\item \textit{[1 pt]} $E(X)=4$, $V(X)=\frac{16}{5}$.

\item \textit{[\num{0,5} pt]} Pour $N$ essais indépendants, la probabilité d'aucun succès est $(1-p)^N$. La probabilité demandée est donc $1-(\frac{4}{5})^N$. Elle est strictement croissante avec $N$, car $0<\frac{4}{5}<1$. Les calculs donnent
\[P_{13}\simeq\num{0,945024}<\num{0,95},\qquad P_{14}\simeq\num{0,956020}\ge\num{0,95}.\]
Ainsi $14$ est le plus petit nombre d'essais convenable.

\item \textit{[\num{0,5} pt]} Notons $n=20$. Pour $j\in\{0,\ldots,n\}$, un résultat ordonné contenant $j$ succès et $n-j$ échecs a, par indépendance, la probabilité $p^j(1-p)^{n-j}$. Il existe $\binom nj$ choix des positions des succès. Ces résultats sont incompatibles, donc
\[P(X=j)=\binom nj p^j(1-p)^{n-j},\qquad p=\frac{1}{5}.\]

\item \textit{[\num{0,5} pt]} Si $B_1,\ldots,B_N$ sont les indicatrices indépendantes des succès, $F_N=\frac{B_1+\cdots+B_N}{N}$. Avec $E(B_i)=p$ et $V(B_i)=p(1-p)$,
\[E(F_N)=\frac{Np}{N}=\frac{1}{5},\qquad V(F_N)=\frac{Np(1-p)}{N^2}=\frac{\frac{4}{25}}{N}.\]
L'indépendance est requise pour l'addition des variances, pas pour la linéarité de l'espérance.

\item \textit{[1 pt]} Bienaymé-Tchebychev appliquée à $F_N$ donne
\[P(|F_N-p|\ge\num{0,02})\le\frac{p(1-p)}{N(\num{0,02})^2}.\]
Pour garantir que cette probabilité est au plus $\num{0,05}$, il suffit de
\[N\ge\frac{\frac{4}{25}}{\num{0,05}\times(\num{0,02})^2}=8000.\]
Par complément, on a alors $P(|F_N-p|<\num{0,02})\ge\num{0,95}$. La taille $8000$ est suffisante ; l'inégalité étant une majoration, elle ne fournit pas nécessairement la taille minimale pour la probabilité exacte.

\end{enumerate}

\fi

\mxnewpage\phantomsection\label{sujet:9-2}\addcontentsline{toc}{chapter}{Blanc final · Zéro 2}

\section*{Blanc final · Zéro 2}

\mxbanner{SUJET ZÉRO · ENTRAÎNEMENT}

\textbf{Durée : 4 h · Note sur 20 · \mxTallied{ex-9-2} exercices indépendants.}\\Calculatrice autorisée. Toute réponse doit être justifiée ; les démarches incomplètes seront prises en compte. Donner les valeurs exactes ; arrondir les probabilités à $10^{-6}$ et les autres approximations à $10^{-4}$, sauf précision différente.

\smallskip\textbf{Nom :}\enspace\dotfill\quad\textbf{Classe :}\enspace\dotfill\quad\textbf{Date :}\enspace\dotfill

\mxtally{ex-9-2}
\subsection*{Exercice 1 · Une fonction logarithmique \hfill 5 points}On considère la fonction $f$ définie sur $\intoo{0}{+\infty}$ par $f(x)=\frac{5+\ln x}x$. On note $\mathcal C_f$ sa courbe représentative.

\begin{enumerate}

\item Déterminer les limites de $f$ en $0^+$ et en $+\infty$. \hfill\textit{(\num{0,5} pt)}

\item Calculer $f' (x)$, puis dresser le tableau de variations de $f$. \hfill\textit{(\num{0,5} pt)}

\item Préciser la valeur du maximum de $f$, puis résoudre $f(x)=0$. \hfill\textit{(\num{0,5} pt)}

\item Étudier la convexité de $f$ et déterminer les coordonnées du point d'inflexion de $\mathcal C_f$. \hfill\textit{(1 pt)}

\item Démontrer que l'équation $f(x)=\frac12 e^{4}$ admet exactement deux solutions dans $\intoo{0}{+\infty}$. \hfill\textit{(\num{0,5} pt)}

\item Déterminer une primitive de $f$ sur $\intoo{0}{+\infty}$. \hfill\textit{(\num{0,5} pt)}

\item Calculer la valeur moyenne de $f$ sur $\intff{1}{e}$. \hfill\textit{(\num{0,5} pt)}

\item Indépendamment des questions précédentes, résoudre $\ln(x-3)+\ln(x+3)=\ln(40)$, après avoir précisé l'ensemble sur lequel cette équation est définie. \hfill\textit{(1 pt)}

\end{enumerate}

\mxtally{ex-9-2}
\subsection*{Exercice 2 · Phénomène périodique \hfill 5 points}Un point se déplace sur une roue tournant à vitesse constante dans un plan vertical. Sa hauteur au-dessus du sol est modélisée par $h(t)=12-4\cos\left(\frac{\pi t}{4}\right)$, où $t$ est exprimé en secondes et $h(t)$ en mètres. On étudie un tour complet, pour $t\in\intff{0}{8}$.

\begin{enumerate}

\item Donner une période du modèle prolongé à $\R$, puis les hauteurs minimale et maximale sur un tour. \hfill\textit{(\num{0,5} pt)}

\item Calculer $h'(t)$ et dresser le tableau de variations de $h$ sur $\intff{0}{8}$. \hfill\textit{(\num{0,5} pt)}

\item Déterminer tous les instants $t\in\intff{0}{8}$ tels que $h(t)=14$. \hfill\textit{(\num{0,5} pt)}

\item Calculer la valeur moyenne de $h$ sur $\intff{0}{8}$. \hfill\textit{(1 pt)}

\item Étudier la convexité de $h$ sur l'intervalle d'étude et déterminer les coordonnées des points d'inflexion de sa courbe. \hfill\textit{(\num{0,5} pt)}

\item Déterminer la durée totale pendant laquelle la hauteur est au moins égale à $14$ m. \hfill\textit{(\num{0,5} pt)}

\item La vitesse verticale est donnée par $h'(t)$. Déterminer ses valeurs maximale et minimale, en mètres par seconde, ainsi que les instants où elles sont atteintes. \hfill\textit{(\num{0,5} pt)}

\item Calculer le déplacement vertical $h(8)-h(0)$, puis la somme de la hauteur parcourue à la montée et de celle parcourue à la descente. Expliquer pourquoi ces deux résultats sont différents. \hfill\textit{(1 pt)}

\end{enumerate}

\mxnewpage\mxcontinued{Blanc final · Zéro 2 · suite}

\mxtally{ex-9-2}
\subsection*{Exercice 3 · Droite, plan et projeté \hfill 5 points}Dans un repère orthonormé de l'espace, on considère le plan $\mathcal P$ d'équation $x+2y+2z=7$ et le point $M\pt{3;0;0}$. La droite $d$ passe par $M$ et admet pour vecteur directeur $\vec n=\pt{1;2;2}$.

\begin{enumerate}

\item Donner une représentation paramétrique de $d$. \hfill\textit{(\num{0,5} pt)}

\item Démontrer que $d$ coupe $\mathcal P$ en un point $H$, puis calculer les coordonnées de $H$. \hfill\textit{(\num{0,5} pt)}

\item Justifier que $H$ est le projeté orthogonal de $M$ sur $\mathcal P$, puis calculer la distance de $M$ à $\mathcal P$. \hfill\textit{(\num{0,5} pt)}

\item On considère le plan $\mathcal Q$ d'équation $y=z$. Démontrer que $\mathcal P\cap\mathcal Q$ est une droite et en donner une représentation paramétrique. \hfill\textit{(1 pt)}

\item Vérifier que $H\in\mathcal Q$. En déduire la distance de $M$ à la droite $\mathcal P\cap\mathcal Q$. \hfill\textit{(\num{0,5} pt)}

\item Déterminer une équation cartésienne du plan parallèle à $\mathcal P$ passant par $M$. \hfill\textit{(\num{0,5} pt)}

\item Déterminer l'angle entre $d$ et la droite $\mathcal P\cap\mathcal Q$. \hfill\textit{(\num{0,5} pt)}

\item On définit le point $N$ par $\vect{MN}=\vec n$. Calculer les distances de $N$ au plan $\mathcal P$ et à la droite $\mathcal P\cap\mathcal Q$, puis les comparer. \hfill\textit{(1 pt)}

\end{enumerate}

\mxtally{ex-9-2}
\subsection*{Exercice 4 · Répétitions et estimation \hfill 5 points}On répète $17$ fois une épreuve à deux issues, appelées succès et échec. Les épreuves sont indépendantes et la probabilité de succès est $p=\frac{3}{10}$ à chaque épreuve. On note $X$ le nombre de succès obtenus.

\begin{enumerate}

\item Justifier que $X$ suit une loi binomiale et préciser ses paramètres. \hfill\textit{(\num{0,5} pt)}

\item Calculer la probabilité d'obtenir exactement deux succès. \hfill\textit{(\num{0,5} pt)}

\item Exprimer à l'aide des événements contraires, puis calculer, la probabilité d'obtenir au moins deux succès. \hfill\textit{(\num{0,5} pt)}

\item Calculer l'espérance et la variance de $X$. \hfill\textit{(1 pt)}

\item On répète maintenant la même épreuve un nombre $N\ge1$ de fois, toujours indépendamment. Déterminer le plus petit entier $N$ pour lequel la probabilité d'obtenir au moins un succès est supérieure ou égale à $\num{0,95}$. Une recherche par tableau ou par boucle est autorisée. \hfill\textit{(\num{0,5} pt)}

\item Pour un entier $j\in\{0,\ldots,17\}$, démontrer la formule de $P(X=j)$ en dénombrant les positions possibles des succès. \hfill\textit{(\num{0,5} pt)}

\item Pour $N$ épreuves indépendantes de même paramètre $p$, on note $F_N$ la fréquence des succès. Démontrer les expressions de $E(F_N)$ et de $V(F_N)$ en fonction de $N$ et de $p$. \hfill\textit{(\num{0,5} pt)}

\item À l'aide de Bienaymé-Tchebychev, déterminer une taille $N$ suffisante pour que $P(|F_N-p|<\num{0,02})\ge\num{0,95}$. La taille obtenue est-elle nécessairement la plus petite possible pour cette probabilité ? Justifier. \hfill\textit{(1 pt)}

\end{enumerate}

\ifmxteacher\mxnewpage

\section*{Corrigé · Blanc final · Zéro 2}\label{cor:9-2}

\subsection*{Exercice 1 · Une fonction logarithmique}\begin{enumerate}

\item \textit{[\num{0,5} pt]} En $0^+$, le numérateur tend vers $-\infty$ et le dénominateur vers $0^+$ : limite $-\infty$ ; en $+\infty$, limite $0$ par croissance comparée.

\item \textit{[\num{0,5} pt]} Par la règle du quotient,
\[f'(x)=\frac{(\frac{1}{x})x-(5+\ln x)}{x^2}=\frac{-4-\ln x}{x^2}.\]
Le dénominateur est positif. Le numérateur s'annule en $s=e^{-4}$, est positif pour $0<x<s$ et négatif pour $x>s$. Ainsi $f$ croît strictement sur $]0,s]$, puis décroît strictement sur $[s,+\infty[$.

\item \textit{[\num{0,5} pt]} Le maximum est atteint uniquement en $s=e^{-4}$ et vaut
\[f(s)=\frac{5- 4}{e^{-4}}=e^{4}.\]
Comme $x>0$, $f(x)=0$ équivaut à $5+\ln x=0$, soit $x=e^{-5}$.

\item \textit{[1 pt]} On dérive $(-4-\ln x)x^{-2}$ :
\[f''(x)=-x^{-3}-2(-4-\ln x)x^{-3}=\frac{2\ln x+7}{x^3}.\]
Le numérateur s'annule en $t=e^{-\frac{7}{2}}$ et passe du signe négatif au signe positif. La fonction est concave sur $]0,t]$, convexe sur $[t,+\infty[$. Elle change de convexité en $t$ et $f(t)=\frac{3}{2t}$ : le point d'inflexion est $I\pt{t;\frac{3}{2t}}$.

\item \textit{[\num{0,5} pt]} Posons $m=\frac{e^{4}}{2}$ et $s=e^{-4}$. La fonction est continue sur $]0,+\infty[$. Sur $]0,s]$, elle est strictement croissante, de limite $-\infty$ en $0$ et de valeur $2m$ en $s$ : le corollaire du TVI donne exactement une solution dans $]0,s[$. Sur $[s,+\infty[$, elle est strictement décroissante de $2m$ vers $0<m$ : exactement une solution dans $]s,+\infty[$. Ces intervalles couvrent le domaine et $f(s)\ne m$ : il y a exactement deux solutions.

\item \textit{[\num{0,5} pt]} Posons $F(x)=5\ln x+\frac12(\ln x)^2$, dérivable pour $x>0$. Alors $F'(x)=\frac{5}{x}+\frac{\ln x}{x}=f(x)$ ; $F$ est une primitive sur $]0,+\infty[$.

\item \textit{[\num{0,5} pt]} Avec cette primitive, $F(e)=\frac{11}{2}$ et $F(1)=0$. L'intégrale vaut donc $\frac{11}{2}$, et la valeur moyenne est
\[\mu=\frac{1}{e-1}\int_1^e f(x)\,\dd x=\frac{11}{2(e-1)}.\]

\item \textit{[1 pt]} Domaine $x>3$ ; $x^2-9=40$, donc $x=7$, la racine négative étant exclue.

\end{enumerate}

\subsection*{Exercice 2 · Phénomène périodique}\begin{enumerate}

\item \textit{[\num{0,5} pt]} Période $8$ s ; minimum $8$ m, maximum $16$ m.

\item \textit{[\num{0,5} pt]} $h'(t)=\pi\sin (\frac{\pi t}{4})$ ; croissance sur $[0,4]$, décroissance sur $[4,8]$.

\item \textit{[\num{0,5} pt]} L'équation équivaut à $\cos (\frac{\pi t}{4})=-\frac{1}{2}$. Sur un tour, $\frac{\pi t}{4}\in[0,2\pi]$ ; les angles convenables sont $\frac{2\pi}{3}$ et $\frac{4\pi}{3}$. Ainsi $t=\frac{8}{3}$ s ou $t=\frac{16}{3}$ s.

\item \textit{[1 pt]} Une primitive est $F(t)=12t-\frac{16}{\pi}\sin (\frac{\pi t}{4})$. Comme $\sin0=\sin(2\pi)=0$,
\[\overline h=\frac{1}{8}\int_0^{8}h(t)\,\dd t=\frac{F(8)-F(0)}{8}=\frac{96}{8}=12\ \mathrm{m}.\]

\item \textit{[\num{0,5} pt]} La dérivée seconde est $h''(t)=\frac{4\pi^2}{16}\cos (\frac{\pi t}{4})$. Son signe est celui du cosinus : positif sur $[0,2[$ et $]6,8]$, négatif entre ces deux intervalles. La fonction est convexe sur $[0,2]$ et $[6,8]$, concave sur $[2,6]$. Les deux changements de convexité donnent les points d'inflexion $\pt{2;12}$ et $\pt{6;12}$.

\item \textit{[\num{0,5} pt]} Entre $\frac{8}{3}$ et $\frac{16}{3}$ ; durée $\frac{8}{3}$ s.

\item \textit{[\num{0,5} pt]} $h'$ atteint $1\pi$ $\mathrm{m\,s^{-1}}$ en $t=2$ et son opposé en $t=6$.

\item \textit{[1 pt]} Déplacement $h(8)-h(0)=0$. Montée de $8$ m puis descente de $8$ m : distance cumulée $16$ m, égale à l'intégrale de $|h'|$.

\end{enumerate}

\subsection*{Exercice 3 · Droite, plan et projeté}\begin{enumerate}

\item \textit{[\num{0,5} pt]} $\pt{x;y;z}=\pt{3+t;2t;2t}$, $t\in\R$.

\item \textit{[\num{0,5} pt]} $3+9t=7$ donne $t=\frac{4}{9}$ ; $H=\pt{\frac{31}{9};\frac{8}{9};\frac{8}{9}}$.

\item \textit{[\num{0,5} pt]} Le vecteur $\vec n=\pt{1;2;2}$ est normal à $\mathcal P$. Comme $H\in\mathcal P$ et $\vect{MH}=\frac{4}{9}\vec n$, $H$ est le projeté orthogonal de $M$ sur le plan. Le repère étant orthonormé, $\|\vec n\|=\sqrt{1+4+4}=3$. Donc
\[d(M,\mathcal P)=MH=|\frac{4}{9}|\,3=\frac{4}{3}.\]

\item \textit{[1 pt]} Poser $y=z=s$ ; $x=7-4s$. Donc $\pt{x;y;z}=\pt{7;0;0}+s\pt{-4;1;1}$.

\item \textit{[\num{0,5} pt]} Oui, ses deux dernières coordonnées sont égales. $H$ est dans la droite et $\vect{MH}$ est orthogonal à toute direction du plan, donc à cette droite. Distance $\frac{4}{3}$.

\item \textit{[\num{0,5} pt]} $x+2y+2z=3$ ; même vecteur normal et passage par $M$.

\item \textit{[\num{0,5} pt]} $\vec n\cdot\pt{-4;1;1}=-4+2+2=0$ : angle $\frac{\pi}{2}$.

\item \textit{[1 pt]} $N=M+\vec n$ est sur la même droite normale que $M$ ; son projeté sur $\mathcal P$ est donc aussi $H$. Or $H\in\mathcal P\cap\mathcal Q$. Pour tout $P$ de cette droite, $NP^2=NH^2+HP^2$, donc sa distance minimale à $N$ est $NH$. Finalement,
\[d(N,\mathcal P)=d(N,\mathcal P\cap\mathcal Q)=\|\vect{NH}\|=|\frac{4}{9}-1|\,\|\vec n\|=\frac{5}{3}.\]

\end{enumerate}

\subsection*{Exercice 4 · Répétitions et estimation}\begin{enumerate}

\item \textit{[\num{0,5} pt]} Répétition d'un nombre fixe $17$ d'épreuves identiques indépendantes à deux issues ; $X\sim\mathcal B\pt{17;\frac{3}{10}}$.

\item \textit{[\num{0,5} pt]} $P(X=2)=\binom{17}2(\frac{3}{10})^2(\frac{7}{10})^{15}\simeq\num{0,058110}$.

\item \textit{[\num{0,5} pt]} $1-P(X=0)-P(X=1)=1-(\frac{7}{10})^{17}-17(\frac{3}{10})(\frac{7}{10})^{16}\simeq\num{0,980725}$.

\item \textit{[1 pt]} $E(X)=\frac{51}{10}$, $V(X)=\frac{357}{100}$.

\item \textit{[\num{0,5} pt]} Pour $N$ essais indépendants, la probabilité d'aucun succès est $(1-p)^N$. La probabilité demandée est donc $1-(\frac{7}{10})^N$. Elle est strictement croissante avec $N$, car $0<\frac{7}{10}<1$. Les calculs donnent
\[P_{8}\simeq\num{0,942352}<\num{0,95},\qquad P_{9}\simeq\num{0,959646}\ge\num{0,95}.\]
Ainsi $9$ est le plus petit nombre d'essais convenable.

\item \textit{[\num{0,5} pt]} Notons $n=17$. Pour $j\in\{0,\ldots,n\}$, un résultat ordonné contenant $j$ succès et $n-j$ échecs a, par indépendance, la probabilité $p^j(1-p)^{n-j}$. Il existe $\binom nj$ choix des positions des succès. Ces résultats sont incompatibles, donc
\[P(X=j)=\binom nj p^j(1-p)^{n-j},\qquad p=\frac{3}{10}.\]

\item \textit{[\num{0,5} pt]} Si $B_1,\ldots,B_N$ sont les indicatrices indépendantes des succès, $F_N=\frac{B_1+\cdots+B_N}{N}$. Avec $E(B_i)=p$ et $V(B_i)=p(1-p)$,
\[E(F_N)=\frac{Np}{N}=\frac{3}{10},\qquad V(F_N)=\frac{Np(1-p)}{N^2}=\frac{\frac{21}{100}}{N}.\]
L'indépendance est requise pour l'addition des variances, pas pour la linéarité de l'espérance.

\item \textit{[1 pt]} Bienaymé-Tchebychev appliquée à $F_N$ donne
\[P(|F_N-p|\ge\num{0,02})\le\frac{p(1-p)}{N(\num{0,02})^2}.\]
Pour garantir que cette probabilité est au plus $\num{0,05}$, il suffit de
\[N\ge\frac{\frac{21}{100}}{\num{0,05}\times(\num{0,02})^2}=10500.\]
Par complément, on a alors $P(|F_N-p|<\num{0,02})\ge\num{0,95}$. La taille $10500$ est suffisante ; l'inégalité étant une majoration, elle ne fournit pas nécessairement la taille minimale pour la probabilité exacte.

\end{enumerate}

\fi

\ifmxteacher

\mxnewpage\phantomsection\label{sujet:9-3}\addcontentsline{toc}{chapter}{Blanc final · A}

\section*{Blanc final · A}

\mxbanner*{SUJET RÉSERVÉ À L'ÉVALUATION}

\textbf{Durée : 4 h · Note sur 20 · \mxTallied{ex-9-3} exercices indépendants.}\\Calculatrice autorisée. Toute réponse doit être justifiée ; les démarches incomplètes seront prises en compte. Donner les valeurs exactes ; arrondir les probabilités à $10^{-6}$ et les autres approximations à $10^{-4}$, sauf précision différente.

\smallskip\textbf{Nom :}\enspace\dotfill\quad\textbf{Classe :}\enspace\dotfill\quad\textbf{Date :}\enspace\dotfill

\mxtally{ex-9-3}
\subsection*{Exercice 1 · Une fonction logarithmique \hfill 5 points}On considère la fonction $f$ définie sur $\intoo{0}{+\infty}$ par $f(x)=\frac{4+\ln x}x$. On note $\mathcal C_f$ sa courbe représentative.

\begin{enumerate}

\item Déterminer les limites de $f$ en $0^+$ et en $+\infty$. \hfill\textit{(\num{0,5} pt)}

\item Calculer $f' (x)$, puis dresser le tableau de variations de $f$. \hfill\textit{(\num{0,5} pt)}

\item Préciser la valeur du maximum de $f$, puis résoudre $f(x)=0$. \hfill\textit{(\num{0,5} pt)}

\item Étudier la convexité de $f$ et déterminer les coordonnées du point d'inflexion de $\mathcal C_f$. \hfill\textit{(1 pt)}

\item Démontrer que l'équation $f(x)=\frac12 e^{3}$ admet exactement deux solutions dans $\intoo{0}{+\infty}$. \hfill\textit{(\num{0,5} pt)}

\item Déterminer une primitive de $f$ sur $\intoo{0}{+\infty}$. \hfill\textit{(\num{0,5} pt)}

\item Calculer la valeur moyenne de $f$ sur $\intff{1}{e}$. \hfill\textit{(\num{0,5} pt)}

\item Indépendamment des questions précédentes, résoudre $\ln(x-1)+\ln(x+1)=\ln(15)$, après avoir précisé l'ensemble sur lequel cette équation est définie. \hfill\textit{(1 pt)}

\end{enumerate}

\mxtally{ex-9-3}
\subsection*{Exercice 2 · Phénomène périodique \hfill 5 points}Un point se déplace sur une roue tournant à vitesse constante dans un plan vertical. Sa hauteur au-dessus du sol est modélisée par $h(t)=7-2\cos\left(\frac{\pi t}{3}\right)$, où $t$ est exprimé en secondes et $h(t)$ en mètres. On étudie un tour complet, pour $t\in\intff{0}{6}$.

\begin{enumerate}

\item Donner une période du modèle prolongé à $\R$, puis les hauteurs minimale et maximale sur un tour. \hfill\textit{(\num{0,5} pt)}

\item Calculer $h'(t)$ et dresser le tableau de variations de $h$ sur $\intff{0}{6}$. \hfill\textit{(\num{0,5} pt)}

\item Déterminer tous les instants $t\in\intff{0}{6}$ tels que $h(t)=8$. \hfill\textit{(\num{0,5} pt)}

\item Calculer la valeur moyenne de $h$ sur $\intff{0}{6}$. \hfill\textit{(1 pt)}

\item Étudier la convexité de $h$ sur l'intervalle d'étude et déterminer les coordonnées des points d'inflexion de sa courbe. \hfill\textit{(\num{0,5} pt)}

\item Déterminer la durée totale pendant laquelle la hauteur est au moins égale à $8$ m. \hfill\textit{(\num{0,5} pt)}

\item La vitesse verticale est donnée par $h'(t)$. Déterminer ses valeurs maximale et minimale, en mètres par seconde, ainsi que les instants où elles sont atteintes. \hfill\textit{(\num{0,5} pt)}

\item Calculer le déplacement vertical $h(6)-h(0)$, puis la somme de la hauteur parcourue à la montée et de celle parcourue à la descente. Expliquer pourquoi ces deux résultats sont différents. \hfill\textit{(1 pt)}

\end{enumerate}

\mxnewpage\mxcontinued{Blanc final · A · suite}

\mxtally{ex-9-3}
\subsection*{Exercice 3 · Droite, plan et projeté \hfill 5 points}Dans un repère orthonormé de l'espace, on considère le plan $\mathcal P$ d'équation $x+2y+2z=8$ et le point $M\pt{1;0;0}$. La droite $d$ passe par $M$ et admet pour vecteur directeur $\vec n=\pt{1;2;2}$.

\begin{enumerate}

\item Donner une représentation paramétrique de $d$. \hfill\textit{(\num{0,5} pt)}

\item Démontrer que $d$ coupe $\mathcal P$ en un point $H$, puis calculer les coordonnées de $H$. \hfill\textit{(\num{0,5} pt)}

\item Justifier que $H$ est le projeté orthogonal de $M$ sur $\mathcal P$, puis calculer la distance de $M$ à $\mathcal P$. \hfill\textit{(\num{0,5} pt)}

\item On considère le plan $\mathcal Q$ d'équation $y=z$. Démontrer que $\mathcal P\cap\mathcal Q$ est une droite et en donner une représentation paramétrique. \hfill\textit{(1 pt)}

\item Vérifier que $H\in\mathcal Q$. En déduire la distance de $M$ à la droite $\mathcal P\cap\mathcal Q$. \hfill\textit{(\num{0,5} pt)}

\item Déterminer une équation cartésienne du plan parallèle à $\mathcal P$ passant par $M$. \hfill\textit{(\num{0,5} pt)}

\item Déterminer l'angle entre $d$ et la droite $\mathcal P\cap\mathcal Q$. \hfill\textit{(\num{0,5} pt)}

\item On définit le point $N$ par $\vect{MN}=\vec n$. Calculer les distances de $N$ au plan $\mathcal P$ et à la droite $\mathcal P\cap\mathcal Q$, puis les comparer. \hfill\textit{(1 pt)}

\end{enumerate}

\mxtally{ex-9-3}
\subsection*{Exercice 4 · Répétitions et estimation \hfill 5 points}On répète $18$ fois une épreuve à deux issues, appelées succès et échec. Les épreuves sont indépendantes et la probabilité de succès est $p=\frac{1}{10}$ à chaque épreuve. On note $X$ le nombre de succès obtenus.

\begin{enumerate}

\item Justifier que $X$ suit une loi binomiale et préciser ses paramètres. \hfill\textit{(\num{0,5} pt)}

\item Calculer la probabilité d'obtenir exactement deux succès. \hfill\textit{(\num{0,5} pt)}

\item Exprimer à l'aide des événements contraires, puis calculer, la probabilité d'obtenir au moins deux succès. \hfill\textit{(\num{0,5} pt)}

\item Calculer l'espérance et la variance de $X$. \hfill\textit{(1 pt)}

\item On répète maintenant la même épreuve un nombre $N\ge1$ de fois, toujours indépendamment. Déterminer le plus petit entier $N$ pour lequel la probabilité d'obtenir au moins un succès est supérieure ou égale à $\num{0,95}$. Une recherche par tableau ou par boucle est autorisée. \hfill\textit{(\num{0,5} pt)}

\item Pour un entier $j\in\{0,\ldots,18\}$, démontrer la formule de $P(X=j)$ en dénombrant les positions possibles des succès. \hfill\textit{(\num{0,5} pt)}

\item Pour $N$ épreuves indépendantes de même paramètre $p$, on note $F_N$ la fréquence des succès. Démontrer les expressions de $E(F_N)$ et de $V(F_N)$ en fonction de $N$ et de $p$. \hfill\textit{(\num{0,5} pt)}

\item À l'aide de Bienaymé-Tchebychev, déterminer une taille $N$ suffisante pour que $P(|F_N-p|<\num{0,02})\ge\num{0,95}$. La taille obtenue est-elle nécessairement la plus petite possible pour cette probabilité ? Justifier. \hfill\textit{(1 pt)}

\end{enumerate}

\ifmxteacher\mxnewpage

\section*{Corrigé · Blanc final · A}\label{cor:9-3}

\subsection*{Exercice 1 · Une fonction logarithmique}\begin{enumerate}

\item \textit{[\num{0,5} pt]} En $0^+$, le numérateur tend vers $-\infty$ et le dénominateur vers $0^+$ : limite $-\infty$ ; en $+\infty$, limite $0$ par croissance comparée.

\item \textit{[\num{0,5} pt]} Par la règle du quotient,
\[f'(x)=\frac{(\frac{1}{x})x-(4+\ln x)}{x^2}=\frac{-3-\ln x}{x^2}.\]
Le dénominateur est positif. Le numérateur s'annule en $s=e^{-3}$, est positif pour $0<x<s$ et négatif pour $x>s$. Ainsi $f$ croît strictement sur $]0,s]$, puis décroît strictement sur $[s,+\infty[$.

\item \textit{[\num{0,5} pt]} Le maximum est atteint uniquement en $s=e^{-3}$ et vaut
\[f(s)=\frac{4- 3}{e^{-3}}=e^{3}.\]
Comme $x>0$, $f(x)=0$ équivaut à $4+\ln x=0$, soit $x=e^{-4}$.

\item \textit{[1 pt]} On dérive $(-3-\ln x)x^{-2}$ :
\[f''(x)=-x^{-3}-2(-3-\ln x)x^{-3}=\frac{2\ln x+5}{x^3}.\]
Le numérateur s'annule en $t=e^{-\frac{5}{2}}$ et passe du signe négatif au signe positif. La fonction est concave sur $]0,t]$, convexe sur $[t,+\infty[$. Elle change de convexité en $t$ et $f(t)=\frac{3}{2t}$ : le point d'inflexion est $I\pt{t;\frac{3}{2t}}$.

\item \textit{[\num{0,5} pt]} Posons $m=\frac{e^{3}}{2}$ et $s=e^{-3}$. La fonction est continue sur $]0,+\infty[$. Sur $]0,s]$, elle est strictement croissante, de limite $-\infty$ en $0$ et de valeur $2m$ en $s$ : le corollaire du TVI donne exactement une solution dans $]0,s[$. Sur $[s,+\infty[$, elle est strictement décroissante de $2m$ vers $0<m$ : exactement une solution dans $]s,+\infty[$. Ces intervalles couvrent le domaine et $f(s)\ne m$ : il y a exactement deux solutions.

\item \textit{[\num{0,5} pt]} Posons $F(x)=4\ln x+\frac12(\ln x)^2$, dérivable pour $x>0$. Alors $F'(x)=\frac{4}{x}+\frac{\ln x}{x}=f(x)$ ; $F$ est une primitive sur $]0,+\infty[$.

\item \textit{[\num{0,5} pt]} Avec cette primitive, $F(e)=\frac{9}{2}$ et $F(1)=0$. L'intégrale vaut donc $\frac{9}{2}$, et la valeur moyenne est
\[\mu=\frac{1}{e-1}\int_1^e f(x)\,\dd x=\frac{9}{2(e-1)}.\]

\item \textit{[1 pt]} Domaine $x>1$ ; $x^2-1=15$, donc $x=4$, la racine négative étant exclue.

\end{enumerate}

\subsection*{Exercice 2 · Phénomène périodique}\begin{enumerate}

\item \textit{[\num{0,5} pt]} Période $6$ s ; minimum $5$ m, maximum $9$ m.

\item \textit{[\num{0,5} pt]} $h'(t)=(\frac{2}{3}\pi)\sin (\frac{\pi t}{3})$ ; croissance sur $[0,3]$, décroissance sur $[3,6]$.

\item \textit{[\num{0,5} pt]} L'équation équivaut à $\cos (\frac{\pi t}{3})=-\frac{1}{2}$. Sur un tour, $\frac{\pi t}{3}\in[0,2\pi]$ ; les angles convenables sont $\frac{2\pi}{3}$ et $\frac{4\pi}{3}$. Ainsi $t=2$ s ou $t=4$ s.

\item \textit{[1 pt]} Une primitive est $F(t)=7t-\frac{6}{\pi}\sin (\frac{\pi t}{3})$. Comme $\sin0=\sin(2\pi)=0$,
\[\overline h=\frac{1}{6}\int_0^{6}h(t)\,\dd t=\frac{F(6)-F(0)}{6}=\frac{42}{6}=7\ \mathrm{m}.\]

\item \textit{[\num{0,5} pt]} La dérivée seconde est $h''(t)=\frac{2\pi^2}{9}\cos (\frac{\pi t}{3})$. Son signe est celui du cosinus : positif sur $[0,\frac{3}{2}[$ et $]\frac{9}{2},6]$, négatif entre ces deux intervalles. La fonction est convexe sur $[0,\frac{3}{2}]$ et $[\frac{9}{2},6]$, concave sur $[\frac{3}{2},\frac{9}{2}]$. Les deux changements de convexité donnent les points d'inflexion $\pt{\frac{3}{2};7}$ et $\pt{\frac{9}{2};7}$.

\item \textit{[\num{0,5} pt]} Entre $2$ et $4$ ; durée $2$ s.

\item \textit{[\num{0,5} pt]} $h'$ atteint $\frac{2}{3}\pi$ $\mathrm{m\,s^{-1}}$ en $t=\frac{3}{2}$ et son opposé en $t=\frac{9}{2}$.

\item \textit{[1 pt]} Déplacement $h(6)-h(0)=0$. Montée de $4$ m puis descente de $4$ m : distance cumulée $8$ m, égale à l'intégrale de $|h'|$.

\end{enumerate}

\subsection*{Exercice 3 · Droite, plan et projeté}\begin{enumerate}

\item \textit{[\num{0,5} pt]} $\pt{x;y;z}=\pt{1+t;2t;2t}$, $t\in\R$.

\item \textit{[\num{0,5} pt]} $1+9t=8$ donne $t=\frac{7}{9}$ ; $H=\pt{\frac{16}{9};\frac{14}{9};\frac{14}{9}}$.

\item \textit{[\num{0,5} pt]} Le vecteur $\vec n=\pt{1;2;2}$ est normal à $\mathcal P$. Comme $H\in\mathcal P$ et $\vect{MH}=\frac{7}{9}\vec n$, $H$ est le projeté orthogonal de $M$ sur le plan. Le repère étant orthonormé, $\|\vec n\|=\sqrt{1+4+4}=3$. Donc
\[d(M,\mathcal P)=MH=|\frac{7}{9}|\,3=\frac{7}{3}.\]

\item \textit{[1 pt]} Poser $y=z=s$ ; $x=8-4s$. Donc $\pt{x;y;z}=\pt{8;0;0}+s\pt{-4;1;1}$.

\item \textit{[\num{0,5} pt]} Oui, ses deux dernières coordonnées sont égales. $H$ est dans la droite et $\vect{MH}$ est orthogonal à toute direction du plan, donc à cette droite. Distance $\frac{7}{3}$.

\item \textit{[\num{0,5} pt]} $x+2y+2z=1$ ; même vecteur normal et passage par $M$.

\item \textit{[\num{0,5} pt]} $\vec n\cdot\pt{-4;1;1}=-4+2+2=0$ : angle $\frac{\pi}{2}$.

\item \textit{[1 pt]} $N=M+\vec n$ est sur la même droite normale que $M$ ; son projeté sur $\mathcal P$ est donc aussi $H$. Or $H\in\mathcal P\cap\mathcal Q$. Pour tout $P$ de cette droite, $NP^2=NH^2+HP^2$, donc sa distance minimale à $N$ est $NH$. Finalement,
\[d(N,\mathcal P)=d(N,\mathcal P\cap\mathcal Q)=\|\vect{NH}\|=|\frac{7}{9}-1|\,\|\vec n\|=\frac{2}{3}.\]

\end{enumerate}

\subsection*{Exercice 4 · Répétitions et estimation}\begin{enumerate}

\item \textit{[\num{0,5} pt]} Répétition d'un nombre fixe $18$ d'épreuves identiques indépendantes à deux issues ; $X\sim\mathcal B\pt{18;\frac{1}{10}}$.

\item \textit{[\num{0,5} pt]} $P(X=2)=\binom{18}2(\frac{1}{10})^2(\frac{9}{10})^{16}\simeq\num{0,283512}$.

\item \textit{[\num{0,5} pt]} $1-P(X=0)-P(X=1)=1-(\frac{9}{10})^{18}-18(\frac{1}{10})(\frac{9}{10})^{17}\simeq\num{0,549716}$.

\item \textit{[1 pt]} $E(X)=\frac{9}{5}$, $V(X)=\frac{81}{50}$.

\item \textit{[\num{0,5} pt]} Pour $N$ essais indépendants, la probabilité d'aucun succès est $(1-p)^N$. La probabilité demandée est donc $1-(\frac{9}{10})^N$. Elle est strictement croissante avec $N$, car $0<\frac{9}{10}<1$. Les calculs donnent
\[P_{28}\simeq\num{0,947665}<\num{0,95},\qquad P_{29}\simeq\num{0,952899}\ge\num{0,95}.\]
Ainsi $29$ est le plus petit nombre d'essais convenable.

\item \textit{[\num{0,5} pt]} Notons $n=18$. Pour $j\in\{0,\ldots,n\}$, un résultat ordonné contenant $j$ succès et $n-j$ échecs a, par indépendance, la probabilité $p^j(1-p)^{n-j}$. Il existe $\binom nj$ choix des positions des succès. Ces résultats sont incompatibles, donc
\[P(X=j)=\binom nj p^j(1-p)^{n-j},\qquad p=\frac{1}{10}.\]

\item \textit{[\num{0,5} pt]} Si $B_1,\ldots,B_N$ sont les indicatrices indépendantes des succès, $F_N=\frac{B_1+\cdots+B_N}{N}$. Avec $E(B_i)=p$ et $V(B_i)=p(1-p)$,
\[E(F_N)=\frac{Np}{N}=\frac{1}{10},\qquad V(F_N)=\frac{Np(1-p)}{N^2}=\frac{\frac{9}{100}}{N}.\]
L'indépendance est requise pour l'addition des variances, pas pour la linéarité de l'espérance.

\item \textit{[1 pt]} Bienaymé-Tchebychev appliquée à $F_N$ donne
\[P(|F_N-p|\ge\num{0,02})\le\frac{p(1-p)}{N(\num{0,02})^2}.\]
Pour garantir que cette probabilité est au plus $\num{0,05}$, il suffit de
\[N\ge\frac{\frac{9}{100}}{\num{0,05}\times(\num{0,02})^2}=4500.\]
Par complément, on a alors $P(|F_N-p|<\num{0,02})\ge\num{0,95}$. La taille $4500$ est suffisante ; l'inégalité étant une majoration, elle ne fournit pas nécessairement la taille minimale pour la probabilité exacte.

\end{enumerate}

\fi

\fi

\ifmxteacher

\mxnewpage\phantomsection\label{sujet:9-4}\addcontentsline{toc}{chapter}{Blanc final · B}

\section*{Blanc final · B}

\mxbanner*{SUJET RÉSERVÉ À L'ÉVALUATION}

\textbf{Durée : 4 h · Note sur 20 · \mxTallied{ex-9-4} exercices indépendants.}\\Calculatrice autorisée. Toute réponse doit être justifiée ; les démarches incomplètes seront prises en compte. Donner les valeurs exactes ; arrondir les probabilités à $10^{-6}$ et les autres approximations à $10^{-4}$, sauf précision différente.

\smallskip\textbf{Nom :}\enspace\dotfill\quad\textbf{Classe :}\enspace\dotfill\quad\textbf{Date :}\enspace\dotfill

\mxtally{ex-9-4}
\subsection*{Exercice 1 · Une fonction logarithmique \hfill 5 points}On considère la fonction $f$ définie sur $\intoo{0}{+\infty}$ par $f(x)=\frac{3+\ln x}x$. On note $\mathcal C_f$ sa courbe représentative.

\begin{enumerate}

\item Déterminer les limites de $f$ en $0^+$ et en $+\infty$. \hfill\textit{(\num{0,5} pt)}

\item Calculer $f' (x)$, puis dresser le tableau de variations de $f$. \hfill\textit{(\num{0,5} pt)}

\item Préciser la valeur du maximum de $f$, puis résoudre $f(x)=0$. \hfill\textit{(\num{0,5} pt)}

\item Étudier la convexité de $f$ et déterminer les coordonnées du point d'inflexion de $\mathcal C_f$. \hfill\textit{(1 pt)}

\item Démontrer que l'équation $f(x)=\frac12 e^{2}$ admet exactement deux solutions dans $\intoo{0}{+\infty}$. \hfill\textit{(\num{0,5} pt)}

\item Déterminer une primitive de $f$ sur $\intoo{0}{+\infty}$. \hfill\textit{(\num{0,5} pt)}

\item Calculer la valeur moyenne de $f$ sur $\intff{1}{e}$. \hfill\textit{(\num{0,5} pt)}

\item Indépendamment des questions précédentes, résoudre $\ln(x-2)+\ln(x+2)=\ln(32)$, après avoir précisé l'ensemble sur lequel cette équation est définie. \hfill\textit{(1 pt)}

\end{enumerate}

\mxtally{ex-9-4}
\subsection*{Exercice 2 · Phénomène périodique \hfill 5 points}Un point se déplace sur une roue tournant à vitesse constante dans un plan vertical. Sa hauteur au-dessus du sol est modélisée par $h(t)=10-3\cos\left(\frac{\pi t}{2}\right)$, où $t$ est exprimé en secondes et $h(t)$ en mètres. On étudie un tour complet, pour $t\in\intff{0}{4}$.

\begin{enumerate}

\item Donner une période du modèle prolongé à $\R$, puis les hauteurs minimale et maximale sur un tour. \hfill\textit{(\num{0,5} pt)}

\item Calculer $h'(t)$ et dresser le tableau de variations de $h$ sur $\intff{0}{4}$. \hfill\textit{(\num{0,5} pt)}

\item Déterminer tous les instants $t\in\intff{0}{4}$ tels que $h(t)=\frac{23}{2}$. \hfill\textit{(\num{0,5} pt)}

\item Calculer la valeur moyenne de $h$ sur $\intff{0}{4}$. \hfill\textit{(1 pt)}

\item Étudier la convexité de $h$ sur l'intervalle d'étude et déterminer les coordonnées des points d'inflexion de sa courbe. \hfill\textit{(\num{0,5} pt)}

\item Déterminer la durée totale pendant laquelle la hauteur est au moins égale à $\frac{23}{2}$ m. \hfill\textit{(\num{0,5} pt)}

\item La vitesse verticale est donnée par $h'(t)$. Déterminer ses valeurs maximale et minimale, en mètres par seconde, ainsi que les instants où elles sont atteintes. \hfill\textit{(\num{0,5} pt)}

\item Calculer le déplacement vertical $h(4)-h(0)$, puis la somme de la hauteur parcourue à la montée et de celle parcourue à la descente. Expliquer pourquoi ces deux résultats sont différents. \hfill\textit{(1 pt)}

\end{enumerate}

\mxnewpage\mxcontinued{Blanc final · B · suite}

\mxtally{ex-9-4}
\subsection*{Exercice 3 · Droite, plan et projeté \hfill 5 points}Dans un repère orthonormé de l'espace, on considère le plan $\mathcal P$ d'équation $x+2y+2z=12$ et le point $M\pt{2;0;0}$. La droite $d$ passe par $M$ et admet pour vecteur directeur $\vec n=\pt{1;2;2}$.

\begin{enumerate}

\item Donner une représentation paramétrique de $d$. \hfill\textit{(\num{0,5} pt)}

\item Démontrer que $d$ coupe $\mathcal P$ en un point $H$, puis calculer les coordonnées de $H$. \hfill\textit{(\num{0,5} pt)}

\item Justifier que $H$ est le projeté orthogonal de $M$ sur $\mathcal P$, puis calculer la distance de $M$ à $\mathcal P$. \hfill\textit{(\num{0,5} pt)}

\item On considère le plan $\mathcal Q$ d'équation $y=z$. Démontrer que $\mathcal P\cap\mathcal Q$ est une droite et en donner une représentation paramétrique. \hfill\textit{(1 pt)}

\item Vérifier que $H\in\mathcal Q$. En déduire la distance de $M$ à la droite $\mathcal P\cap\mathcal Q$. \hfill\textit{(\num{0,5} pt)}

\item Déterminer une équation cartésienne du plan parallèle à $\mathcal P$ passant par $M$. \hfill\textit{(\num{0,5} pt)}

\item Déterminer l'angle entre $d$ et la droite $\mathcal P\cap\mathcal Q$. \hfill\textit{(\num{0,5} pt)}

\item On définit le point $N$ par $\vect{MN}=\vec n$. Calculer les distances de $N$ au plan $\mathcal P$ et à la droite $\mathcal P\cap\mathcal Q$, puis les comparer. \hfill\textit{(1 pt)}

\end{enumerate}

\mxtally{ex-9-4}
\subsection*{Exercice 4 · Répétitions et estimation \hfill 5 points}On répète $20$ fois une épreuve à deux issues, appelées succès et échec. Les épreuves sont indépendantes et la probabilité de succès est $p=\frac{3}{10}$ à chaque épreuve. On note $X$ le nombre de succès obtenus.

\begin{enumerate}

\item Justifier que $X$ suit une loi binomiale et préciser ses paramètres. \hfill\textit{(\num{0,5} pt)}

\item Calculer la probabilité d'obtenir exactement deux succès. \hfill\textit{(\num{0,5} pt)}

\item Exprimer à l'aide des événements contraires, puis calculer, la probabilité d'obtenir au moins deux succès. \hfill\textit{(\num{0,5} pt)}

\item Calculer l'espérance et la variance de $X$. \hfill\textit{(1 pt)}

\item On répète maintenant la même épreuve un nombre $N\ge1$ de fois, toujours indépendamment. Déterminer le plus petit entier $N$ pour lequel la probabilité d'obtenir au moins un succès est supérieure ou égale à $\num{0,95}$. Une recherche par tableau ou par boucle est autorisée. \hfill\textit{(\num{0,5} pt)}

\item Pour un entier $j\in\{0,\ldots,20\}$, démontrer la formule de $P(X=j)$ en dénombrant les positions possibles des succès. \hfill\textit{(\num{0,5} pt)}

\item Pour $N$ épreuves indépendantes de même paramètre $p$, on note $F_N$ la fréquence des succès. Démontrer les expressions de $E(F_N)$ et de $V(F_N)$ en fonction de $N$ et de $p$. \hfill\textit{(\num{0,5} pt)}

\item À l'aide de Bienaymé-Tchebychev, déterminer une taille $N$ suffisante pour que $P(|F_N-p|<\num{0,02})\ge\num{0,95}$. La taille obtenue est-elle nécessairement la plus petite possible pour cette probabilité ? Justifier. \hfill\textit{(1 pt)}

\end{enumerate}

\ifmxteacher\mxnewpage

\section*{Corrigé · Blanc final · B}\label{cor:9-4}

\subsection*{Exercice 1 · Une fonction logarithmique}\begin{enumerate}

\item \textit{[\num{0,5} pt]} En $0^+$, le numérateur tend vers $-\infty$ et le dénominateur vers $0^+$ : limite $-\infty$ ; en $+\infty$, limite $0$ par croissance comparée.

\item \textit{[\num{0,5} pt]} Par la règle du quotient,
\[f'(x)=\frac{(\frac{1}{x})x-(3+\ln x)}{x^2}=\frac{-2-\ln x}{x^2}.\]
Le dénominateur est positif. Le numérateur s'annule en $s=e^{-2}$, est positif pour $0<x<s$ et négatif pour $x>s$. Ainsi $f$ croît strictement sur $]0,s]$, puis décroît strictement sur $[s,+\infty[$.

\item \textit{[\num{0,5} pt]} Le maximum est atteint uniquement en $s=e^{-2}$ et vaut
\[f(s)=\frac{3- 2}{e^{-2}}=e^{2}.\]
Comme $x>0$, $f(x)=0$ équivaut à $3+\ln x=0$, soit $x=e^{-3}$.

\item \textit{[1 pt]} On dérive $(-2-\ln x)x^{-2}$ :
\[f''(x)=-x^{-3}-2(-2-\ln x)x^{-3}=\frac{2\ln x+3}{x^3}.\]
Le numérateur s'annule en $t=e^{-\frac{3}{2}}$ et passe du signe négatif au signe positif. La fonction est concave sur $]0,t]$, convexe sur $[t,+\infty[$. Elle change de convexité en $t$ et $f(t)=\frac{3}{2t}$ : le point d'inflexion est $I\pt{t;\frac{3}{2t}}$.

\item \textit{[\num{0,5} pt]} Posons $m=\frac{e^{2}}{2}$ et $s=e^{-2}$. La fonction est continue sur $]0,+\infty[$. Sur $]0,s]$, elle est strictement croissante, de limite $-\infty$ en $0$ et de valeur $2m$ en $s$ : le corollaire du TVI donne exactement une solution dans $]0,s[$. Sur $[s,+\infty[$, elle est strictement décroissante de $2m$ vers $0<m$ : exactement une solution dans $]s,+\infty[$. Ces intervalles couvrent le domaine et $f(s)\ne m$ : il y a exactement deux solutions.

\item \textit{[\num{0,5} pt]} Posons $F(x)=3\ln x+\frac12(\ln x)^2$, dérivable pour $x>0$. Alors $F'(x)=\frac{3}{x}+\frac{\ln x}{x}=f(x)$ ; $F$ est une primitive sur $]0,+\infty[$.

\item \textit{[\num{0,5} pt]} Avec cette primitive, $F(e)=\frac{7}{2}$ et $F(1)=0$. L'intégrale vaut donc $\frac{7}{2}$, et la valeur moyenne est
\[\mu=\frac{1}{e-1}\int_1^e f(x)\,\dd x=\frac{7}{2(e-1)}.\]

\item \textit{[1 pt]} Domaine $x>2$ ; $x^2-4=32$, donc $x=6$, la racine négative étant exclue.

\end{enumerate}

\subsection*{Exercice 2 · Phénomène périodique}\begin{enumerate}

\item \textit{[\num{0,5} pt]} Période $4$ s ; minimum $7$ m, maximum $13$ m.

\item \textit{[\num{0,5} pt]} $h'(t)=(\frac{3}{2}\pi)\sin (\frac{\pi t}{2})$ ; croissance sur $[0,2]$, décroissance sur $[2,4]$.

\item \textit{[\num{0,5} pt]} L'équation équivaut à $\cos (\frac{\pi t}{2})=-\frac{1}{2}$. Sur un tour, $\frac{\pi t}{2}\in[0,2\pi]$ ; les angles convenables sont $\frac{2\pi}{3}$ et $\frac{4\pi}{3}$. Ainsi $t=\frac{4}{3}$ s ou $t=\frac{8}{3}$ s.

\item \textit{[1 pt]} Une primitive est $F(t)=10t-\frac{6}{\pi}\sin (\frac{\pi t}{2})$. Comme $\sin0=\sin(2\pi)=0$,
\[\overline h=\frac{1}{4}\int_0^{4}h(t)\,\dd t=\frac{F(4)-F(0)}{4}=\frac{40}{4}=10\ \mathrm{m}.\]

\item \textit{[\num{0,5} pt]} La dérivée seconde est $h''(t)=\frac{3\pi^2}{4}\cos (\frac{\pi t}{2})$. Son signe est celui du cosinus : positif sur $[0,1[$ et $]3,4]$, négatif entre ces deux intervalles. La fonction est convexe sur $[0,1]$ et $[3,4]$, concave sur $[1,3]$. Les deux changements de convexité donnent les points d'inflexion $\pt{1;10}$ et $\pt{3;10}$.

\item \textit{[\num{0,5} pt]} Entre $\frac{4}{3}$ et $\frac{8}{3}$ ; durée $\frac{4}{3}$ s.

\item \textit{[\num{0,5} pt]} $h'$ atteint $\frac{3}{2}\pi$ $\mathrm{m\,s^{-1}}$ en $t=1$ et son opposé en $t=3$.

\item \textit{[1 pt]} Déplacement $h(4)-h(0)=0$. Montée de $6$ m puis descente de $6$ m : distance cumulée $12$ m, égale à l'intégrale de $|h'|$.

\end{enumerate}

\subsection*{Exercice 3 · Droite, plan et projeté}\begin{enumerate}

\item \textit{[\num{0,5} pt]} $\pt{x;y;z}=\pt{2+t;2t;2t}$, $t\in\R$.

\item \textit{[\num{0,5} pt]} $2+9t=12$ donne $t=\frac{10}{9}$ ; $H=\pt{\frac{28}{9};\frac{20}{9};\frac{20}{9}}$.

\item \textit{[\num{0,5} pt]} Le vecteur $\vec n=\pt{1;2;2}$ est normal à $\mathcal P$. Comme $H\in\mathcal P$ et $\vect{MH}=\frac{10}{9}\vec n$, $H$ est le projeté orthogonal de $M$ sur le plan. Le repère étant orthonormé, $\|\vec n\|=\sqrt{1+4+4}=3$. Donc
\[d(M,\mathcal P)=MH=|\frac{10}{9}|\,3=\frac{10}{3}.\]

\item \textit{[1 pt]} Poser $y=z=s$ ; $x=12-4s$. Donc $\pt{x;y;z}=\pt{12;0;0}+s\pt{-4;1;1}$.

\item \textit{[\num{0,5} pt]} Oui, ses deux dernières coordonnées sont égales. $H$ est dans la droite et $\vect{MH}$ est orthogonal à toute direction du plan, donc à cette droite. Distance $\frac{10}{3}$.

\item \textit{[\num{0,5} pt]} $x+2y+2z=2$ ; même vecteur normal et passage par $M$.

\item \textit{[\num{0,5} pt]} $\vec n\cdot\pt{-4;1;1}=-4+2+2=0$ : angle $\frac{\pi}{2}$.

\item \textit{[1 pt]} $N=M+\vec n$ est sur la même droite normale que $M$ ; son projeté sur $\mathcal P$ est donc aussi $H$. Or $H\in\mathcal P\cap\mathcal Q$. Pour tout $P$ de cette droite, $NP^2=NH^2+HP^2$, donc sa distance minimale à $N$ est $NH$. Finalement,
\[d(N,\mathcal P)=d(N,\mathcal P\cap\mathcal Q)=\|\vect{NH}\|=|\frac{10}{9}-1|\,\|\vec n\|=\frac{1}{3}.\]

\end{enumerate}

\subsection*{Exercice 4 · Répétitions et estimation}\begin{enumerate}

\item \textit{[\num{0,5} pt]} Répétition d'un nombre fixe $20$ d'épreuves identiques indépendantes à deux issues ; $X\sim\mathcal B\pt{20;\frac{3}{10}}$.

\item \textit{[\num{0,5} pt]} $P(X=2)=\binom{20}2(\frac{3}{10})^2(\frac{7}{10})^{18}\simeq\num{0,027846}$.

\item \textit{[\num{0,5} pt]} $1-P(X=0)-P(X=1)=1-(\frac{7}{10})^{20}-20(\frac{3}{10})(\frac{7}{10})^{19}\simeq\num{0,992363}$.

\item \textit{[1 pt]} $E(X)=6$, $V(X)=\frac{21}{5}$.

\item \textit{[\num{0,5} pt]} Pour $N$ essais indépendants, la probabilité d'aucun succès est $(1-p)^N$. La probabilité demandée est donc $1-(\frac{7}{10})^N$. Elle est strictement croissante avec $N$, car $0<\frac{7}{10}<1$. Les calculs donnent
\[P_{8}\simeq\num{0,942352}<\num{0,95},\qquad P_{9}\simeq\num{0,959646}\ge\num{0,95}.\]
Ainsi $9$ est le plus petit nombre d'essais convenable.

\item \textit{[\num{0,5} pt]} Notons $n=20$. Pour $j\in\{0,\ldots,n\}$, un résultat ordonné contenant $j$ succès et $n-j$ échecs a, par indépendance, la probabilité $p^j(1-p)^{n-j}$. Il existe $\binom nj$ choix des positions des succès. Ces résultats sont incompatibles, donc
\[P(X=j)=\binom nj p^j(1-p)^{n-j},\qquad p=\frac{3}{10}.\]

\item \textit{[\num{0,5} pt]} Si $B_1,\ldots,B_N$ sont les indicatrices indépendantes des succès, $F_N=\frac{B_1+\cdots+B_N}{N}$. Avec $E(B_i)=p$ et $V(B_i)=p(1-p)$,
\[E(F_N)=\frac{Np}{N}=\frac{3}{10},\qquad V(F_N)=\frac{Np(1-p)}{N^2}=\frac{\frac{21}{100}}{N}.\]
L'indépendance est requise pour l'addition des variances, pas pour la linéarité de l'espérance.

\item \textit{[1 pt]} Bienaymé-Tchebychev appliquée à $F_N$ donne
\[P(|F_N-p|\ge\num{0,02})\le\frac{p(1-p)}{N(\num{0,02})^2}.\]
Pour garantir que cette probabilité est au plus $\num{0,05}$, il suffit de
\[N\ge\frac{\frac{21}{100}}{\num{0,05}\times(\num{0,02})^2}=10500.\]
Par complément, on a alors $P(|F_N-p|<\num{0,02})\ge\num{0,95}$. La taille $10500$ est suffisante ; l'inégalité étant une majoration, elle ne fournit pas nécessairement la taille minimale pour la probabilité exacte.

\end{enumerate}

\fi

\fi

